\section{Alcance final del problema}

El proyecto resolverá tres cosas para los siniestros de enero a septiembre de 2026: la caracterización de la gravedad entre grupos definidos con las circunstancias del parte (OE1); la evaluación de la clasificación retrospectiva con al menos cinco tipos de modelos de aprendizaje automático clásico frente a dos referencias sencillas, con los conjuntos~A y~B (OE2); y la exposición del modelo seleccionado en una interfaz de consulta de carácter académico (OE3).

Quedan fuera, por estar relacionados pero exceder lo que los datos y el plazo permiten: anticipar siniestros o emitir alertas previas; establecer causas o responsabilidades; estudiar la gravedad de las lesiones y los fallecimientos posteriores al evento; incorporar factores externos como el clima o el fenómeno de El Niño; usar datos personales; comparar con otros años; emplear aprendizaje profundo; y construir un sistema operativo para las autoridades. La exclusión del clima se decidió en la Tarea~2 porque, con un solo periodo y sin información territorial suficientemente detallada, no habría una comparación confiable; queda como trabajo futuro.

El límite temporal merece una precisión adicional. Con enero--septiembre de 2026 es posible revisar la estabilidad entre ventanas del mismo año, pero no sostener afirmaciones sobre dos semestres completos ni sobre variabilidad anual. Cualquier conclusión que se presente con esos alcances sería incompatible con R-05 y con el criterio CA-08.
